\documentclass[a4paper]{article}

\date{}

\usepackage{amssymb}

\begin{document}
\setcounter{page}{1}

\title{Simultaneous Extensions of States on Quantum Logics}

\author{
Dominika Bure\v{s}ov\'{a}*, Pavel Pt\'{a}k, Jan \v{S}evic\\
Czech Technical University in Prague\\[1ex]
\texttt{* buresdo2@fel.cvut.cz}}


\maketitle

\begin{abstract}
   Suppose that $L$ is a quantum logic (= an orthomodular poset) and
$A,B$ are Boolean subalgebras of $L$. Suppose further that
$s_1(\mathrm{resp.}\; s_2)$ is a state on $A$
(resp. on $B$). We ask when both $s_1$ and $s_2$
allow for a simultaneous extension over $L$.
Under a natural necessity condition on
$A,B,s_1,s_2$ we find a sufficient condition
for the existence of the extension.
This condition is the requirement of an abundance
of two-valued states on $L$. We then show
by the Greechie pasting technique that this condition is far from necessary even for state-space-rich logics.
\end{abstract}

\begin{thebibliography}{9}

\bibitem{RaoRao}
M.~M.~Rao and Z.~D.~Rao.
\newblock \emph{Theory of Charges.}
\newblock Academic Press, 1983.

\bibitem{PtakPulmannova}
P.~Pt\'ak and S.~Pulmannov\'a.
\newblock \emph{Orthomodular Structures as Quantum Logics.}
\newblock Kluwer Academic Publishers, 1991.

\end{thebibliography}

\section*{Acknowledgements}
This research was supported by the GACR (project No. 25-20013L) and the SGS (project No. SGS26/073/OHK3/1T/13).


\end{document}



